\section{Experimental Protocols and Reproducibility}
\label{sec:supp_protocols}


\subsection{Datasets, Labels, and Prompts}
\label{sec:supp_data_prompts}

\begin{table}[!ht]
\centering
\caption{\textbf{Exact language interfaces.}
Templates and grade strings are reproduced verbatim.
Semicolons separate grade strings and are not inserted into the prompts.}
\label{tab:supp_f_prompts}
\begingroup
\fontsize{8}{9.5}\selectfont
\setlength{\tabcolsep}{3pt}
\renewcommand{\arraystretch}{1}
\begin{tabularx}{\columnwidth}{
  @{}
  >{\raggedright\arraybackslash}p{0.21\columnwidth}
  >{\raggedright\arraybackslash}X
  @{}
}
\toprule
Task & Language interface\\
\midrule
Adience & \texttt{a photo of a person aged \{grade\} years.} & \texttt{0 to 2}, \texttt{4 to 6}, \texttt{8 to 12}, \texttt{15 to 20}, \texttt{25 to 32}, \texttt{38 to 43}, \texttt{48 to 53}, \texttt{60 to 100} \\
HCI & \texttt{a historical color photograph taken in the \{grade\}.} & \texttt{1930s}, \texttt{1940s}, \texttt{1950s}, \texttt{1960s}, \texttt{1970s} \\
FairFace & \texttt{a photo of a person aged \{grade\} years.} & \texttt{0 to 2}, \texttt{3 to 9}, \texttt{10 to 19}, \texttt{20 to 29}, \texttt{30 to 39}, \texttt{40 to 49}, \texttt{50 to 59}, \texttt{60 to 69}, \texttt{more than 70} \\
KonIQ & \texttt{a photograph with \{grade\} image quality.} & \texttt{bad}, \texttt{poor}, \texttt{fair}, \texttt{good}, \texttt{excellent} \\
SCUT-FBP5500 & \texttt{a photo of a face with beauty score \{grade\} out of five.} & \texttt{one}, \texttt{two}, \texttt{three}, \texttt{four}, \texttt{five} \\
DDR & \texttt{a retinal fundus photograph showing \{grade\}.} & \texttt{no diabetic retinopathy}, \texttt{mild non-proliferative diabetic retinopathy}, \texttt{moderate non-proliferative diabetic retinopathy}, \texttt{severe non-proliferative diabetic retinopathy}, \texttt{proliferative diabetic retinopathy} \\
\bottomrule

\end{tabularx}
\endgroup
\end{table}

\begin{table*}[!t]
\centering
\caption{\textbf{Nine adapter configurations.}
Trainable components, score constructions, objectives, and optimization.
CE denotes cross-entropy; SORD uses the soft targets defined in Sec.~\ref{sec:supp_adapters}.
All unspecified base CLIP parameters remain frozen.}
\label{tab:supp_f_adapters}
\begingroup
\fontsize{8}{9.5}\selectfont
\setlength{\tabcolsep}{4pt}
\renewcommand{\arraystretch}{1}
\begin{tabularx}{\textwidth}{
  @{}
  >{\raggedright\arraybackslash}p{0.15\textwidth}
  >{\raggedright\arraybackslash}p{0.36\textwidth}
  >{\raggedright\arraybackslash}X
  @{}
}
\toprule
Configuration & Components, scores, and objective & Optimization and selection\\
\midrule
CoOp-style &
Shared four-token context initialized with \texttt{a photo of a}; fixed lexical ranks.
$S_k=100f^\top p_k$; CE. &
SGD, learning rate 0.002, momentum 0.9, weight decay 0.
Cosine schedule for 100 epochs.
Select minimum NLL at epochs 50 and 100; earlier checkpoint on ties.\\
\addlinespace[2pt]

OrdinalCLIP-style &
Learned context and three base rank tokens, one token per base rank; linear rank interpolation~\cite{li2022ordinalclip}.
$S_k=100f^\top p_k$.
CE plus text-to-image KL divergence with equal weights. &
Adam, learning rate $10^{-4}$, weight decay 0, 50 epochs.
Multiply learning rate by 0.1 at epoch 30 for age tasks, and at epochs 25 and 35 for HCI/KonIQ.
Select minimum NLL over epochs; earlier checkpoint on ties.\\
\addlinespace[2pt]

KgCoOp-style &
CoOp components and score.
Objective:
$\mathrm{CE}+\frac{8}{K}\sum_{k=1}^{K}\|p_k-p_k^0\|_2^2$,
using all $K$ normalized text features. &
Same SGD, cosine schedule, checkpoint candidates, and selection rule as CoOp-style.\\
\addlinespace[2pt]

TaskRes-style (CE) &
Only a zero-initialized $K\times512$ prototype residual $\Delta W$ is trained.
$S=100\operatorname{norm}(P+0.5\Delta W)f$.
CE. &
Adam with learning rate in $\{0.0002,0.002\}$, $\epsilon=10^{-8}$, weight decay 0.
Cosine schedule, maximum 60 epochs; first-epoch learning rate $10^{-5}$; patience 10.
Select minimum NLL across epochs and learning rates, retaining the earlier candidate on exact ties.\\
\addlinespace[2pt]

TaskRes-style (SORD) &
Same prototype residual and score as TaskRes-style (CE); SORD supervision. &
Same optimization and selection as TaskRes-style (CE).\\
\addlinespace[2pt]

Tip-Adapter-F (CE) &
Unconstrained support keys $W$, initialized with unit image features; fixed one-hot values $V$.
$S=T+\alpha V^\top\exp[-\beta(\mathbf1-Wf)]$.
The exponential is elementwise; CE. &
AdamW with learning rate in $\{10^{-4},10^{-3}\}$, $\epsilon=10^{-4}$, weight decay 0.01.
Cosine schedule, maximum 60 epochs; patience 10.
Fit at $\alpha=1$, $\beta=5.5$.
Retain the minimum-NLL epoch per learning rate, then select learning rate,
$\alpha\in\{0.1,1,10\}$, and
$\beta\in\{1,2,5.5\}$ jointly by NLL.
Exact ties follow ascending learning-rate/$\alpha$/$\beta$ order.\\
\addlinespace[2pt]

Tip-Adapter-F (SORD) &
Same keys, values, and score as Tip-Adapter-F (CE); SORD supervision. &
Same optimization, checkpoint selection, and parameter grids as Tip-Adapter-F (CE).\\
\addlinespace[2pt]

M3 &
Rank-8 LoRA on the final visual attention QKV and output weights, plus a $512\to128\to K$ GELU residual head.
$S_k=100z_\theta^\top p_k^0+h_k$.
The head input is centered and scaled using the native training mean and RMS.
CE.
$T$ uses the unadapted image encoder. &
AdamW with visual learning rate $10^{-4}$ and head learning rate $10^{-3}$; weight decay $10^{-3}$.
Maximum 60 epochs; patience 10.
Select minimum NLL; earlier checkpoint on ties.\\
\addlinespace[2pt]

TEXT\_ONLY &
An independent visual LoRA fit with the same visual construction as M3.
The residual head is disabled and frozen; $S_k=100z_\theta^\top p_k^0$, distinct from frozen $T$; CE. &
AdamW, learning rate $10^{-4}$, weight decay $10^{-3}$.
Maximum 60 epochs; patience 10.
Select minimum NLL; earlier checkpoint on ties.\\
\bottomrule

\end{tabularx}

\medskip
\begin{minipage}{\textwidth}
\noindent\textit{Prompt initialization.}
CoOp-style and KgCoOp-style use the suffixes \texttt{person aged \{grade\} years} for age, \texttt{\{grade\}} for HCI, and \texttt{\{grade\} image quality} for KonIQ.
OrdinalCLIP-style uses the task-specific initialization strings ``age estimation: the age of the person is'' for age tasks, ``historical dating: a photo taken in the'' for HCI, and ``aesthetics assessment: the quality of the photo is'' for KonIQ.
\end{minipage}
\endgroup
\end{table*}

\begin{table}[!t]
\centering
\caption{\textbf{Target-grade support allocations.}
(a) Pre-mask quotas and per-target counts at 0/10/20\% support.
Fixed-total has 256 images on all four tasks.
(b) Fixed-seen totals on Adience and HCI.}
\label{tab:supp_f_support}
\begingroup
\fontsize{8}{9.5}\selectfont
\setlength{\tabcolsep}{3pt}
\renewcommand{\arraystretch}{1}
\textbf{(a) Pool quotas and target counts}\par
\smallskip
\begin{tabularx}{\columnwidth}{
  @{}
  >{\raggedright\arraybackslash}p{0.22\columnwidth}
  >{\raggedright\arraybackslash}X
  >{\centering\arraybackslash}p{0.22\columnwidth}
  @{}
}
\toprule
Task & Pool quota & 0/10/20\%\\
\midrule
Adience & 32 per grade & 0/4/7\\
HCI / KonIQ & 52 for grade 1; 51 otherwise & 0/6/11\\
FairFace & 29 for grades 1--4; 28 otherwise & 0/3/6\\
\bottomrule
\end{tabularx}

\medskip
\textbf{(b) Fixed-seen sample totals}\par
\smallskip
\begin{tabularx}{\columnwidth}{
  @{}
  >{\raggedright\arraybackslash}X
  r
  >{\centering\arraybackslash}p{0.36\columnwidth}
  @{}
}
\toprule
Task and target pattern & Seen & Total: 0/10/20\%\\
\midrule
Adience, one grade & 224 & 224/228/231\\
Adience, two grades & 192 & 192/200/206\\
HCI, low endpoint & 204 & 204/210/215\\
HCI, middle or high & 205 & 205/211/216\\
HCI, adjacent or separated pair & 154 & 154/166/176\\
HCI, both endpoints & 153 & 153/165/175\\
\bottomrule
\end{tabularx}

\endgroup
\end{table}

\begin{table}[!t]
\centering
\caption{\textbf{Validation-only selection.}
All selectors use the bias grid and error budget in Sec.~\ref{sec:supp_calibration_protocol}.
Blend strengths are $\lambda\in\{0,0.1,\ldots,1.5,2\}$.}
\label{tab:supp_f_selectors}
\begingroup
\fontsize{8}{9.5}\selectfont
\setlength{\tabcolsep}{3pt}
\renewcommand{\arraystretch}{1}
\begin{tabularx}{\columnwidth}{
  @{}
  >{\raggedright\arraybackslash}p{0.19\columnwidth}
  >{\raggedright\arraybackslash}X
  @{}
}
\toprule
Selector & Sequence and choice\\
\midrule
Class bias & For each condition let $q_i(Z)=\max_{k\in\mathcal O}Z_k(x_i)-\max_{k\in\mathcal U}Z_k(x_i)$ on seen validation. Using $Z\in\{S,S^{\mathrm{TIE}}\}$, set $c_{\min}=\lfloor\min q\rfloor-1$, $c_{\max}=\lceil\max q\rceil+1$. Candidates $\{-\infty\}\cup\{c_{\min},c_{\min}+1,\ldots,c_{\max}\}$. Largest feasible $c$ under uncalibrated Adapted nMAE$+0.005$ (numerical tolerance $10^{-15}$); each path chooses its own. $-\infty$ suppresses missing grades; empty feasible set is blocked. Full-$K$ argmax ties choose the lowest grade. \\
Blend-S & $\lambda\in\{0,0.1,\ldots,1.5,2\}$. First minimize uncalibrated seen-validation nMAE of $T+\lambda(S-T)$ within the same risk budget; largest $\lambda$ among exact minimum-risk ties. Then select its own largest feasible bias on the condition's frozen shared grid above. No feasible bias: blocked, no replacement offset. \\
Blend-A & Same $17$ strengths and condition-specific bias grid. Independently select the largest feasible bias for every $\lambda$; discard infeasible candidates. Minimize resulting calibrated seen-validation nMAE; largest $\lambda$ among exact equal-risk ties. If all candidates are infeasible, block the path. Test scores do not select either deployment strength. \\
Binary control & Its own unit-step grid uses the same floor/ceiling rule applied to its validation score switches alone, plus $-\infty$. Same uncalibrated Adapted budget and largest-feasible rule; a failed guard remains blocked. \\
\bottomrule

\end{tabularx}
\endgroup
\end{table}

The main benchmarks are Adience~\cite{eidinger2014age} ($K=8$), HCI~\cite{palermo2012dating} ($K=5$), FairFace~\cite{karkkainen2021fairface} ($K=9$), and KonIQ-10k~\cite{hosu2020koniq} ($K=5$).
Each image has one ordered grade and is evaluated against the full $K$-class label space.
Test sizes per seed are 4,377, 198--199, 14,356, and 2,010, respectively.

KonIQ grades are training-split MOS quintiles with cut points 2.7016, 3.1385, 3.4167, and 3.6400.
Table~\ref{tab:supp_f_prompts} gives the exact reference templates and ordered grade strings.
Learned-context and ordinal-prompt configurations retain the same grade identities as their zero-shot references.


\subsection{Adapter Configurations}
\label{sec:supp_adapters}

All nine main configurations use OpenAI CLIP ViT-B/16~\cite{radford2021clip} and produce aligned pre-softmax scores for the same grade vocabulary.
They cover learned context prompts~\cite{zhou2022coop}, ordinal rank prompts~\cite{li2022ordinalclip}, knowledge-preserving prompts~\cite{yao2023kgcoop}, prototype residuals~\cite{yu2023taskres}, and trainable cache adapters~\cite{zhang2022tipadapter}.


The CE and SORD configurations differ in their supervision targets.
SORD uses distance-based soft ordinal labels~\cite{diaz2019sord}, here
$q_j\propto\exp[-(j-y)^2/2]$ over all $K$ grades.
M3 combines visual low-rank adaptation (LoRA)~\cite{hu2022lora} with a residual head.
TEXT\_ONLY trains the visual LoRA independently with the residual head disabled and text embeddings frozen.
Its name refers to scoring with the frozen text prototypes alone, without M3's residual head.

In Table~\ref{tab:supp_f_adapters}, $f$ is the unit frozen-image feature, $p_k$ and $p_k^0$ are adapted and reference unit text features, and $P$ stacks reference text features as rows.
The adapted visual feature is $z_\theta$, and $h_k$ is the residual-head output for grade $k$.
The operator $\operatorname{norm}$ normalizes matrix rows.
For cache adapters, $W$ contains support keys and $V$ their one-hot grade values.

All configurations use batch size 32.
Checkpoint and hyperparameter selection use hard-label negative log-likelihood (NLL) on seen-grade validation images.
The main KgCoOp-style recipe differs from the formula-matched control in Sec.~\ref{sec:supp_ordinal_adaptation}.


\subsection{Missing Masks and Support Budgets}
\label{sec:supp_f_masks}

The six missing-grade masks are listed in Table~\ref{tab:supp_missing_masks}.
The main protocol uses seeds 101, 202, and 303, with 256 training and 64 validation images drawn only from seen grades.
Counts are balanced across the $K-|\mathcal U|$ seen grades, differing by at most one image.
Every test split retains all $K$ grades.
The main matrix therefore contains $9\times4\times6\times3=648$ conditions.

The support experiments start from a same-seed, balanced, 256-image training-only pool.
Let $n_k^{\mathrm{pool}}$ be the pre-mask quota of target grade $k$.
At support fraction $r$, its sample count is
$n_k(0)=0$ and
$n_k(r)=\max(1,\lceil r n_k^{\mathrm{pool}}\rceil)$ for $r>0$.
Thus, 10\% and 20\% refer to the target grade's pool quota.

Fixed-total keeps 256 images by removing seen-grade images in round-robin order as target images are added.
Its target and seen counts are
$\sum_{k\in\mathcal U}n_k(r)$ and
$256-\sum_{k\in\mathcal U}n_k(r)$.
Fixed-seen retains the exact seen-image IDs from the masked pool and adds target images, using its own matched 0\% baseline.
Both protocols use the six masks and three seeds above.
Table~\ref{tab:supp_f_support} gives all sample allocations.


\subsection{Scores, Calibration, and Blend Selection}
\label{sec:supp_calibration_protocol}

All score constructions use aligned grade indices and a common pre-softmax scale.
The displacement $D=S-T$ is computed before calibration.
Missing-grade calibration adds a class bias, following the calibrated-stacking formulation~\cite{chao2016gzsl}.
For scores $Z$, prediction uses
$\arg\max_k\{Z_k(x)+c_Z^*\mathbf1[k\in\mathcal U]\}$.

The validation budget uses normalized mean absolute error,
\begin{equation}
\operatorname{nMAE}_{\mathrm{seen}}(Z)
=
\frac{1}{N_{\mathrm{val}}(K-1)}
\sum_{i\in\mathrm{val}}
|\hat y_Z(x_i)-y_i|,
\label{eq:supp_f_nmae}
\end{equation}
where $N_{\mathrm{val}}$ is the number of seen-grade validation images.
A bias is feasible when
$\operatorname{nMAE}_{\mathrm{seen}}(Z+c\mathbf1_{\mathcal U})
\le
\operatorname{nMAE}_{\mathrm{seen}}(S)+0.005$,
with numerical tolerance $10^{-15}$.

For validation image $i$, define
$v_i(Z)=\max_{k\in\mathcal O}Z_k(x_i)-\max_{k\in\mathcal U}Z_k(x_i)$.
The shared condition-specific grid uses
$\mathcal V=\{v_i(S),v_i(S^{\mathrm{TIE}}):i\in\mathrm{val}\}$,
with
$c_{\min}=\lfloor\min\mathcal V\rfloor-1$,
$c_{\max}=\lceil\max\mathcal V\rceil+1$, and
$\mathcal C=\{-\infty\}\cup\{c_{\min},c_{\min}+1,\ldots,c_{\max}\}$.
Each path selects its own largest feasible bias.
An empty feasible set is reported as infeasible.
Full-$K$ argmax ties select the lowest grade.

Table~\ref{tab:supp_f_selectors} specifies the two Blend selectors.
Both use seen validation only.
The best Blend in Sec.~\ref{sec:supp_best_blend} is instead the retrospective test maximum on the 17-point grid.
FullSup retains the same total 256/64 training/validation budgets, distributed across all $K$ grades.


\subsection{Metrics, Aggregation, and Uncertainty}
\label{sec:supp_f_metrics}

For each adjacent missing--seen pair $(m,n)$, pair AUROC ranks images by $Z_m-Z_n$, using test images of those two true grades, with $m$ positive and ties counting one half.
MS-AUC is the arithmetic mean over these pairs.
It is computed from finite pre-calibration scores; the $-\infty$ candidate is used only to suppress missing-grade predictions.

Missing MAE first averages absolute grade error within each missing grade, then across missing grades.
Seen and Overall MAE average over their respective test images.
All MAEs are in grade units.

For the main adapter matrix, results first average over six masks and three seeds within each adapter--dataset setting, then equally over the nine configurations.
The interior analysis uses three masks.
Extension results use the masks, seeds, or folds stated in Sec.~\ref{sec:supp_generality}.

FullSup gap closure is computed within each adapter--dataset setting before aggregation.
A ratio is reported when the FullSup improvement over Adapted exceeds $10^{-8}$ and its original paired 95\% interval lies above zero.
The only excluded MS-AUC ratio is TaskRes-style (SORD) on Adience.

The ranking intervals in Fig.~\ref{fig:supp_config_curves} and Table~\ref{tab:supp_best_blend} use 10{,}000 paired group-bootstrap replicates with fixed checkpoints.
Both methods use the same resampled test groups.
Adience groups follow the supplied \texttt{user\_id} metadata; HCI, FairFace, and KonIQ use content-proxy groups.
These intervals are pointwise.

For DDR, content-proxy groups are used because patient identifiers are unavailable.

The counts 34/36 and 35/36 denote positive point estimates across adapter--dataset settings for all masks and interior masks, respectively, at the original strength.
The count 59/64 denotes nonzero dataset--strength comparisons whose paired pointwise 95\% intervals lie above zero.

\subsection{Figure Construction and Implementation}
\label{sec:supp_f_figures}

The main paper's Fig.~1 uses KonIQ, CoOp-style adaptation, seed 101, and a missing middle grade.
Grade 3 (\emph{fair}) is compared with seen grade 2 (\emph{poor}), using 411 and 401 test images.
Pair AUROC is computed from the original per-image score margins.
The displayed offsets are $c^*=6$ for Adapted and $c^*=0$ for TIE.

The main paper's Fig.~2 uses 81 interior-mask conditions per dataset and the 17-point strength grid, with paired bootstrap intervals for TIE minus Blend.
The margin-shift illustration in Fig.~1b is schematic.
The score-processing algorithm and its $O(K)$ cost are given in Sec.~\ref{sec:supp_inference}.
